%%% OBS! Detta är inte en officiell examensmall %%%
%%% OBS! Detta är inte en officiell examensmall %%%
%
% Orginalrepo av Anders Linden:  https://github.com/Anders-Linden/JTH-rapportmall-examensarbete
%
% Uppdaterad av David Isaksson 2018:
% https://github.com/Granddave/JTH-rapportmall-examensarbete
%
% Baserad på följande docx (Hämtad 2018-02-13)
% http://ju.se/download/18.448e5a1414f64cc6e183ea5e/1441623803036/Mall+f%C3%B6r+examensarbete.dotx
% Hämtad från http://ju.se/student/studier/examensarbete-uppsats.html
%
%%% OBS! Detta är inte en officiell examensmall %%%
%%% OBS! Detta är inte en officiell examensmall %%%


\documentclass[12pt,a4paper]{article}
% Article duger till allt som inte har kapitel.
% LaTeX justerar radlängden automatiskt för att det ska bli
% lätt att läsa, vilket kan leda till allt för stora marginaler om man
% inte väljer teckenstorlek 12pt. Alternativt kan man ha dubbla
% kolumner.

\usepackage[margin=34mm]{geometry}
% Ger något smalare marginaler än standard geometrin på article klassen (ca 36mm).

%\usepackage{showframe}
% För att se en ram över dokumentet

\usepackage{blindtext}
% \blindtext genererar lite test-text


\usepackage{lmodern}

\usepackage[utf8]{inputenc}
\usepackage[T1]{fontenc}
% För riktiga å, ä och ö:
% Den första säger att dokumentet är kodat som UTF-8.
% Den andra säger att vi ska använda fonter som är speciellt anpassade
% för Västeuropeiska språk, bland annat kommer prickarna över "ö" lite
% närmare o:et.
% Det räcker ofta, men inte alltid, att bara ange den senare. Det bör
% räcka att bara ange den första, men blir inte _fullt_ så snyggt.

%\usepackage{pslatex}
% Kommentera bort för Times new roman-typsnitt

\usepackage[hidelinks]{hyperref}
% Klickbara länkar

%\usepackage[swedish,english]{babel}
\usepackage[swedish]{babel}
% Svensk avstavning (om den är installerad) och svenska rubriker på
% sådant som LaTeX lägger till: ,
% innehållsförteckning...

\usepackage{parskip}

\usepackage{verbatim}
% För comment-environments

\newenvironment{poliabstract}[1]
  {\renewcommand{\abstractname}{#1}\begin{abstract}}
  {\end{abstract}}
% Eftersom rapporten skall innehålla engelsk sammanfattning också.

\usepackage{amsmath}
% AMS är the American Mathematical Society. De har egna
% LaTeX-kommandon. De är bra, och jag har ingen koll på vilka
% kommandon jag använder som kommer från LaTeX och vilka som kommer
% från AMS

\usepackage{listings}
% tillåter att formatera kod

\lstset{literate=%
    {Ö}{{\"O}}1
    {Ø}{{\O}}1
    {Ä}{{\"A}}1
    {Å}{{\AA}}1
    {Ü}{{\"U}}1
    {ß}{{\ss}}2
    {ü}{{\"u}}1
    %{ä}{{\"a}}1
    %{å}{{\aa}}1
    %{ö}{{\"o}}1
    {ø}{{\o}}1
    {æ}{{\ae}}1
    {Æ}{{\AE}}1
}
% Fixar så att åäö och andra tecken fungerar att använda i kodfiler

\usepackage{setspace, caption}
\DeclareCaptionLabelSeparator*{spaced}{\\[2ex]}
\captionsetup[table]{
    labelfont={it, bf},
    textfont=it,
%    format=plain,
    justification=justified,
%    singlelinecheck=false,
    labelsep=period,%spaced,
    width=.95\textwidth
%    skip=0pt
    }
\captionsetup[figure]{
    labelsep=period,
    labelfont={it, bf},
    textfont=it,
    justification=justified,
    singlelinecheck=false,
    width=.95\textwidth
    %font=doublespacing
    }
\setlength{\belowcaptionskip}{-5pt}


\renewcommand\appendix{\par
  \setcounter{section}{0}
  \setcounter{subsection}{0}
  \setcounter{figure}{0}
  \setcounter{table}{0}
  \renewcommand\thesection{Bilaga \Alph{section}}
  \renewcommand\thefigure{\Alph{section}\arabic{figure}}
}
% Lägger till rubriker för Bilagorna


\usepackage{units}
% Enheter ska skrivas i upprätt stil med seriffer, och separeras från
% föregående siffror med ett tunt mellanslag - tunnare än ett
% vanligt. Använd i texten/matteläget som: \unit[1,34]{m}

\usepackage{icomma}
% I icke-engelsk skrift används decimalkomma, i engelsk är kommatecken
% i matteläge enbart en koordinatavskiljare. Därför sätter LaTeX i
% vanliga fall automatiskt in ett extra mellanslag efter alla
% kommatecken i matteläge. Detta kommando ger ett intelligent
% kommatecken som förstår om det är decimalavskiljare eller
% koordinatavskiljare, beroende på om det står något mellanslag
% efter. Använd aldrig decimalpunkt i icke-engelsk skrift.

\usepackage{tabularx, booktabs}
% Tabellformatering

\usepackage{color}
% Färglägg text med ex. \textcolor{red}{Min röda text}

\usepackage{graphicx}
% Dessa behövs när man inkluderar figurer från xfig, i
% latex+postscript format (brukar bli snyggast). Använder man vanliga
% eps-figurer behöver man inte color-paketet. Det finns både paketet
% "graphics" och "graphicx", med något olika syntax. Om alla bilder
% har rätt storlek när man importerar dem behöver man inte bry sig om
% skillnaderna.

\usepackage{subcaption}
% Captions för subfigurer

\usepackage{bbm}
\newcommand{\N}{\ensuremath{\mathbbm{N}}}
\newcommand{\Z}{\ensuremath{\mathbbm{Z}}}
\newcommand{\Q}{\ensuremath{\mathbbm{Q}}}
\newcommand{\R}{\ensuremath{\mathbbm{R}}}
\newcommand{\C}{\ensuremath{\mathbbm{C}}}
% Här hittar jag på nya kommandon för att beteckna mängden av de
% naturliga talen, heltalen, de rationella, de reella och de
% komplexa. Vill man ha en annan font kan man i stället använda
% \mathbb{N} och så vidare, men måste då inkludera följande rad:

\usepackage{amssymb}
\graphicspath{}
\newcommand{\rd}{\ensuremath{\mathrm{d}}}
\newcommand{\id}{\ensuremath{\,\rd}}
% \rd står för rakt d, \id står för integral-d. Anledningen är att man
% ska kunna skilja på d som en variabel (diameter, avstånd) och d som
% en operator. I derivator och integraler ska d:et stå upprätt, och i
% integraler också med ett litet mellanrum mellan d och föregående
% uttryck.

\setlength{\parindent}{0pt}
% Tar bort indentering vid ny paragraf

\usepackage{gensymb}
% Används för att ge tillgång till grader- (°) symbolen.

\usepackage{float}
%används för korrekt figurplacering

\usepackage{fancyhdr}
% Detta paket gör att det går att header och footer med, vilket krävs för att hålla sig efter den officiella mallen. 

\usepackage{makeidx}
% Index för uppslagsord, vill du lägga till ett ord i saklistan måste du använda dig av \index{} med ordet innanför "måsvingarna".
% Exempel: Flygande bäckasiner \index{bäckasiner} söka hwila på mjuka tuvor.


\let\stdsection\section
\renewcommand\section{\newpage\stdsection}
% För att göra texen mer lättläst och avdelad så läggs alltid en ny huvudrubrik (section) på en helt ny sida.


\newcommand{\blankpage}{
    \newpage
    \thispagestyle{empty}
    \mbox{}
    \newpage
}
% Vissa sidor skall vara tomma eftersom att rapporten kommer att vara utskriven dubbelsidigt. för att skapa en helt tom sida anropas det med \blankpage 
% Det medför en helt tom sida som inte har någon header eller footer.

\usepackage{rotating}
\usepackage[yyyymmdd]{datetime}
\renewcommand{\dateseparator}{-}
\usepackage[absolute]{textpos}
% Varje gång du genererar ett dokument kommer dagens datum att tas fram


%%
%% Referenser
%%
\usepackage[
  backend = biber,
  style   = apa,
  natbib  = true,
  %url     = true,
  sorting = none
]{biblatex}
\addbibresource{referenser}
% Möjlighet att skapa refernslista, med hjälp av BibTex-formatet.

\DeclareLanguageMapping{swedish}{swedish-apa}
\DefineBibliographyStrings{swedish}{andothers = {et al.}}



\renewcommand{\UrlFont}{\tt}
% Url:er får monospace font

\makeindex
% Innan dokumentet börjar måste vi starta igång makeindex så den har möjlighet att söka igenom dokumentet på sakord innan saklistan genereras i slutet av dokumentet.

%%%%%%%%%%%%%%%%%%%%%%%%%%%
%% Dokumentet börjar här %%
%%%%%%%%%%%%%%%%%%%%%%%%%%%
\begin{document}
\pagestyle{empty}

\includegraphics[width=7 cm]{Bilder/JU_logga_stor.jpg}
\vfill

% Titeln på ditt examensarbete
\textbf{\Huge \textcolor{red}{Titel på examensarbetet \\ på två rader} }
\\[0.6cm]

% Undertitel
\textbf{\large \textcolor{red}{Undertitel på en eller\\flera rader} }
\vfill


\textbf{Huvudområde:} \textit{\textcolor{red}{Ämne} } \\
\textbf{Författare:} \textit{\textcolor{red}{Förnamn Efternamn, Förnamn Efternamn} }\\
\textbf{Handledare:} \textit{\textcolor{red}{Förnamn Efternamn} }\\
\textbf{Jönköping} \the\year{} \monthname


%%%% Sida två %%%%
\newpage
\pagestyle{fancy}


\vspace*{\fill}

Detta examensarbete är utfört vid Tekniska Högskolan i Jönköping inom
[se huvudområde på föregående sida].
Författarna svarar själva för framförda åsikter, slutsatser och resultat. \\[0.5em]
\begin{tabular}{@{}ll}
Examinator: & \textcolor{red}{Förnamn Efternamn} \\
Handledare: & \textcolor{red}{Förnamn Efternamn} \\
Omfattning: & \textcolor{red}{7,5 hp eller 9 hp (grundnivå), 15 hp (grundnivå), 30 hp (avancerad nivå)} \\
Datum:      & \today
\end{tabular}

\let\Sectionmark\sectionmark
\def\sectionmark#1{\def\Sectionname{#1}\Sectionmark{#1}}
% Skriver ut överstanivån av sektion på varje sida

\renewcommand{\headrulewidth}{0pt}
\renewcommand{\footrulewidth}{0.4pt}
\lfoot{Postadress:\\ Box 1026 \\ 551 11 Jönköping}
\cfoot{Besöksadress:\\ Gjuterigatan 5}
\rfoot{Telefon: \\ 036-10 10 00 (vx)}

\pagebreak




\renewcommand{\headrulewidth}{0.4pt}

\lfoot{}
\cfoot{\thepage}
\rfoot{}
\pagenumbering{roman}

%\fancyfoot[C]{}
%%\chead{Abstract}
\chead{\textit{Abstract} }
%\selectlanguage{english}
%\begin{poliabstract}{Abstract}
\textbf{\large{Abstract}}\\[0.3cm]
\noindent\textbf{Purpose –} \textcolor{red}{XXX XXX XXX XXX XXX XXX XXX XXX XXX XXX XXX XXX XXX XXX XXX XXX XXX XXX XXX XXX}

\noindent\textbf{Method -} \textcolor{red}{XXX XXX XXX XXX XXX XXX XXX XXX XXX XXX XXX XXX XXX XXX XXX XXX XXX XXX XXX XXX}

\noindent\textbf{Findings -} \textcolor{red}{XXX XXX XXX XXX XXX XXX XXX XXX XXX XXX XXX XXX XXX XXX XXX XXX XXX XXX XXX XXX}

\noindent\textbf{Implications -} \textcolor{red}{XXX XXX XXX XXX XXX XXX XXX XXX XXX XXX XXX XXX XXX XXX XXX XXX XXX XXX XXX XXX}

\noindent\textbf{Limitations -} \textcolor{red}{XXX XXX XXX XXX XXX XXX XXX XXX XXX XXX XXX XXX XXX XXX XXX XXX XXX XXX XXX XXX}

\noindent\textbf{Keywords -} \textcolor{red}{XXX XXX XXX XXX XXX XXX XXX XXX XXX XXX XXX XXX XXX XXX XXX XXX XXX XXX XXX XXX}

%\end{poliabstract}

\newpage%\blankpage

%%\chead{Sammanfattning}
\chead{\textit{Sammanfattning}}
\selectlanguage{swedish}
%\begin{poliabstract}{Sammanfattning}
\textbf{\large{Sammanfattning}}\\[0.3cm]
\noindent\textbf{Syfte –} \textcolor{red}{XXX XXX XXX XXX XXX XXX XXX XXX XXX XXX XXX XXX XXX XXX XXX XXX XXX XXX XXX XXX}

\noindent\textbf{Metod -} \textcolor{red}{XXX XXX XXX XXX XXX XXX XXX XXX XXX XXX XXX XXX XXX XXX XXX XXX XXX XXX XXX XXX}

\noindent\textbf{Resultat -} \textcolor{red}{XXX XXX XXX XXX XXX XXX XXX XXX XXX XXX XXX XXX XXX XXX XXX XXX XXX XXX XXX XXX}

\noindent\textbf{Implikationer -} \textcolor{red}{XXX XXX XXX XXX XXX XXX XXX XXX XXX XXX XXX XXX XXX XXX XXX XXX XXX XXX XXX XXX}

\noindent\textbf{Begränsningar -} \textcolor{red}{XXX XXX XXX XXX XXX XXX XXX XXX XXX XXX XXX XXX XXX XXX XXX XXX XXX XXX XXX XXX}

\noindent\textbf{Nyckelord -} \textcolor{red}{XXX XXX XXX XXX XXX XXX XXX XXX XXX XXX XXX XXX XXX XXX XXX XXX XXX XXX XXX XXX}

\textcolor{red}{Sammanfattningen är en kort beskrivning (minirapport) av rapportens innehåll och bör omfatta högst en A4-sida. Syftet med sammanfattningen är att läsaren snabbt ska få en uppfattning om arbetets förutsättningar och resultat. Det är därför viktigt att sammanfattningen innehåller så konkreta uppgifter som möjligt. Det är viktigt att du lägger mycket tid på att formulera dig eftersom sammanfattningen ofta är det som avgör om rapporten bedöms som värd att läsa eller inte.}
%\end{poliabstract}

\newpage

\lhead{ }
\chead{ \textit {Innehållsförteckning}}
\rhead{ }
\renewcommand{\contentsname}{Innehållsförteckning}
\tableofcontents
\newpage

\newpage

\fancyhead[L]{}
\fancyhead[C]{\textit{\Sectionname}}
\fancyhead[R]{}
\fancyfoot[C]{\thepage}
\pagenumbering{arabic}
\setcounter{page}{1}
\input{inledning.tex}

\section{Teoretisk bakgrund}
Efter att du gjort en genomgång av aktuell litteratur inom det område som ditt examensarbete fokuserar väljer du ut relevanta källor/teorier kopplade till ditt syfte och dina frågeställningar och beskriver det i detta avsnitt. Det är viktigt att du refererar till samtliga källor som du använder i detta kapitel enligt reglerna för referenshantering som nämndes tidigare. Observera att texten inte får vara avskriven från en artikel eller bok utan skall vara av sammanfattande karaktär. Avskrift får endast göras vid ”citat” och bör användas sparsamt.\\ \\ 
För dig som arbetar med design av produkter/system kan detta avsnitt inbegripa olika teoretiska koncept och designteorier, designfilosofier \index{designfilosofier} mm som inspirerat dig i din utveckling av den design du utarbetar. \\ \\
Kontrollera att din teoretiska bakgrund har en tydlig koppling till dina frågeställningar.
\\ \\
Nedan följer ett exempel på hur referenser fungerar.
\\ \\
Lorem ipsum dolor sit amet, consectetur adipiscing elit. Nulla sed rutrum orci. Vestibulum ante ipsum primis in faucibus orci luctus et ultrices posuere cubilia Curae; Vestibulum nec arcu nisl. In diam erat, sodales sit amet tincidunt vel, rutrum a risus \citep{latexcompanion}. Nullam tristique lectus a lectus rhoncus ac fringilla est aliquam. Duis eu dui in urna ultrices pulvinar. Integer leo nibh, pretium sed vehicula eget, bibendum a dui. Proin urna orci, ornare non malesuada ac, gravida sit amet orci \citet{einstein}. Fusce consectetur tellus ut risus dictum pulvinar. Duis turpis dolor, accumsan quis facilisis congue, fringilla a turpis. Aenean vestibulum, nisl vel lobortis facilisis, nulla nibh ornare sem, quis bibendum elit elit tempus arcu. Proin libero dui, porttitor sed sodales sed, luctus vel massa \citep{knuthwebsite}. Integer pellentesque metus vel purus tristique lacinia. Nullam magna orci, mollis at facilisis vitae, sollicitudin sit amet velit. In hac habitasse platea dictumst \citet{latexcompanion,knuthwebsite}. 


\input{metod_genomforande.tex}

\input{resultat.tex}

\section{Diskussion och slutsatser}
I detta kapitel knyter du ihop säcken.


\subsection{I resultatdiskusion / Diskusion av designprocessen}
Här diskuterar du dina analyserade resultat och utvärderar dem utifrån ditt syfte och dina frågeställningar. Här är det meningen att du skall reflektera över dina resultat, varför du fick just dessa resultat, hur de förhåller de sig till vad andra kommit fram till jämfört med din teoribakgrund.\index{teoribakgrund}\\ \\
Du som gör ett designarbete diskuterar här din egen designprocess (eftersom den är en viktig del av ”resultatet” i ett examensarbete) och hur slutprodukten blev i förhållande till vad du tänkt dig enligt syfte och frågeställningar.\\ \\
Ett bra tips för att få en bra struktur på detta avsnitt och som samtidigt hjälper dig att hålla en röd tråd i rapporen är att du inleder med att repetera ditt syfte och sedan utgår från dina frågeställningar som rubriker (det kan du även göra i föregående kapitel om du vill). Då underlättar du för både dig själv och läsaren – det hjälper dig att se till att du verkligen besvarar dina frågeställningar och uppfyller syftet med examensarbetet och det blir lätt för läsaren att följa hur du gjort detta.

\subsection{Metoddiskussion}
Du inleder med en Metoddiskussion under egen rubrik. Här diskuterar du dina metodval och genomförandet i stort och de styrkor/svagheter som du upplever finns i detta avseende. Diskutera hur väl du tycker att du lyckades åstadkomma det du avsåg utifrån ditt syfte och dina frågeställningar genom dina metodval och tillvägagångssätt. Vad fungerade bra/mindre bra? Vad kunde du ha gjort bättre? Hur väl har du uppfyllt kraven på validitet och reliabilitet?

\subsection{Slutsatser och rekommendationer}
Under rubriken Slutsatser summerar du kortfattat de viktigaste huvudpunkterna.\\ \\
Avslutningsvis bör du ge förslag och rekommendationer till hur arbetet skulle kunna fortsättas/utvecklas ytterligare.

\newpage
\input{terminologi.tex}

\newpage

\addcontentsline{toc}{section}{Referenser}
\lhead{ }
\chead{ \textit {Referenser}}
\rhead{ }

\printbibliography

\printindex
\addcontentsline{toc}{section}{Sakregister}


\appendix
\lhead{ }
\chead{ \textit {Bilagor}}
\rhead{ }
\input{bilagor.tex}


%\clearpage\thispagestyle{empty}\null
% Tom sistasida

\end{document}
